\section{Ordered pairs from independent digits}
\label{sec:couplings}

The rational-translation identities of Lemma~\ref{lem:rational-graphs}
will now be transferred to pairs with many opportunities to cross the
thresholds in \eqref{eq:ratio-crossing}. If both endpoint laws were Haar,
coordinate order would make each nonnegative difference $u_i(y)-u_i(x)$
have mean zero, leaving no threshold crossings. We therefore allow a
controlled change in the endpoint marginals. What matters is a uniform
square-integrable bound for their marginal
densities and a product structure that survives conditioning on either
endpoint. We first isolate the density estimate, which also applies to
the independent-coordinate construction in the appendix.

\begin{lemma}[Densities of independent coordinates]
\label{lem:product-density-bound}
Let $(\Omega_i,\mathcal F_i,\mu_i)$, $i\ge1$, be probability spaces and
let $h_i\ge0$ satisfy $\int h_i\,d\mu_i=1$. If
\[
 \sum_{i\ge1}\|h_i-1\|_{L^2(\mu_i)}^2\le C<\infty,
\]
then the product probability measure
$\nu=\bigotimes_{i\ge1}(h_i\mu_i)$ has a density with respect to
$\mu=\bigotimes_{i\ge1}\mu_i$ whose $L^2(\mu)$ norm is at most
$e^{C/2}$.

In particular, suppose two probability densities $h_i^0,h_i^1$ at each coordinate
satisfy
\[
 \|h_i^a-1\|_{L^2(\mu_i)}^2\le b_i\quad(a=0,1),
 \qquad \sum_{i\ge1}b_i\le C.
\]
For every deterministic sequence $\epsilon_i\in\{0,1\}$, the product
law with coordinate densities $h_i^{\epsilon_i}$ has the same bound.
\end{lemma}

\begin{proof}
The mean-one assumption gives
$\|h_i\|_2^2=1+\|h_i-1\|_2^2$. Hence each finite product density obeys
\[
 \left\|\prod_{i=1}^n h_i\right\|_{L^2(\mu)}^2
 =\prod_{i=1}^n\bigl(1+\|h_i-1\|_2^2\bigr)
 \le \exp\left(\sum_{i=1}^n\|h_i-1\|_2^2\right)\le e^C.
\]
For a bounded cylinder function $\varphi$ depending on the first $n$
coordinates, integration under $\nu$ is therefore bounded in absolute
value by $e^{C/2}\|\varphi\|_{L^2(\mu)}$. These cylinder functions are
dense in $L^2(\mu)$, so the integration functional extends continuously
to that space. The Riesz representation theorem gives a real
$H\in L^2(\mu)$ with norm at most $e^{C/2}$ representing the extension.
The finite signed measure $H\mu$ and the probability measure $\nu$
agree on all cylinder sets. Uniqueness of finite measures on the
product sigma-algebra gives $H\mu=\nu$; thus $H\ge0$ almost everywhere
and $\int H\,d\mu=1$. Applying the first assertion to
$h_i=h_i^{\epsilon_i}$ proves the hybrid assertion.
\end{proof}

Write $r_i^*=q_{i+1}/q_i$. A Haar-uniform point of $X$ has the expansion
\begin{equation}\label{eq:digit-expansion}
 x=\frac{A_0}{q_1}+\sum_{i\ge1}\frac{A_i}{q_{i+1}},
 \qquad
 A_0\in\{0,\ldots,q_1-1\},\quad
 A_i\in\{0,\ldots,r_i^*-1\},
\end{equation}
with independent uniform digits. Indeed, specifying a finite prefix selects
an interval whose length is the reciprocal of the product of the relevant
radices. Conversely these intervals determine the Haar measure. The
ambiguity at endpoints of such intervals is a null set. In particular,
$A_i=\lfloor r_i^*u_i(x)\rfloor$, and $u_i(x)$ is the tail of
\eqref{eq:digit-expansion} starting with digit $A_i$, after multiplication
by $q_i$. We couple the independent digits; each character coordinate
$u_i$ depends on an overlapping tail of these digits.

Fix a block index $j$. We construct a probability law $\mathbb P_j$ on a
pair $(x,y)$, retaining auxiliary labels. Use the same uniform initial
digit for both endpoints. Most coordinate pairs will have equal raw
values. Rare top and bottom components force crossings even after one
endpoint has been fixed, and an extra component in block $j$ places
the endpoints close to opposite ends of the unit interval. These
components give tail attenuation while keeping the marginal densities
uniformly square integrable. Independently at every $i\in J_s$, choose a
raw pair $(X_i,Y_i)$ by the following mixture, where
\[
 \rho_s(w)=\kappa_s\sqrt{\frac{w}{m_s}}.
\]
For each of the two labelled widths $w=a_s,b_s$, include these two
components, each of probability $\rho_s(w)$:
\begin{align*}
 \text{top:}\quad &X_i\text{ uniform on }[0,1-w/2),
 &Y_i\text{ uniform on }[1-w/2,1),\\
 \text{bottom:}\quad &X_i\text{ uniform on }[0,w/2),
 &Y_i\text{ uniform on }[w/2,1).
\end{align*}
The two choices within each displayed component are independent. Keep
the width labels distinct even if their numerical values coincide. If
$s=j$, also include a \emph{tagged} component of probability $1/m_s$, with
independent uniforms
\[
 X_i\in[0,1/(2m_s)),\qquad
 Y_i\in[1-1/(2m_s),1).
\]
With the remaining probability set $X_i=Y_i$ to a uniform point of
$[0,1)$. These probabilities have a positive remainder: the parameter
bounds give $\rho_s(w)/w\le\kappa_s$, $w\le1/4$, and $m_s\ge200$.
Define the endpoint digits by
\[
 A_i=\lfloor r_i^*X_i\rfloor,
 \qquad B_i=\lfloor r_i^*Y_i\rfloor,
\]
and use \eqref{eq:digit-expansion} to form $x$ from the $A_i$ and $y$
from the $B_i$. Thus $X_i,Y_i$ choose the first digit of the corresponding
circle coordinates; they are not themselves those coordinates.

\begin{lemma}[Uniform marginal bounds]\label{lem:digit-marginals}
There is a finite constant $H_0$, independent of $j$, such that the
marginal laws of $x$ and $y$ have densities with respect to $\mu$ of
$L^2(\mu)$ norm at most $H_0$. The same assertion holds for an endpoint
formed by any deterministic choice of the $x$-digit or the $y$-digit at
each coordinate, retaining the common uniform initial digit.
\end{lemma}

\begin{proof}
Either raw marginal density is at most $10$. For each width its top and
bottom contributions together are at most $4\rho_s(w)/w$, the tag
contributes at most $2$, and the diagonal component contributes at most
$1$.

We also need a sharper integral estimate. A uniform density on an
interval of length at least $w/2$ differs from $1$ by a mean-zero
function of squared $L^2$ norm at most $2/w$. For a tagged interval the
corresponding bound is $2m_s$. Express a raw marginal density minus $1$
as the sum of its at most five weighted mean-zero perturbations. The
finite-sum Cauchy--Schwarz inequality gives
\begin{equation}\label{eq:one-digit-square}
 \int_0^1 d_i(t)^2\,dt
 \le 1+50\left(\frac{\kappa_s^2}{m_s}
          +\frac{\mathbf 1_{\{s=j\}}}{m_s}\right),
 \qquad i\in J_s.
\end{equation}
In fact the four width terms contribute at most
$40\kappa_s^2/m_s$, and the tag at most $10/m_s$. Passing from raw
variables to digits replaces $d_i$ by its average on each radix bin, so
Jensen's inequality preserves \eqref{eq:one-digit-square}.

Use the uniform law at each digit as the reference probability measure.
The pairs of digits are independent across coordinates, and each
endpoint digit density has mean one. Since $|J_s|=m_s$, summing
\eqref{eq:one-digit-square} minus one over the coordinates gives the
common bound $50(\sum_s\kappa_s^2+1)$, for either endpoint and every
deterministic choice between them. The initial digit has density one.
Lemma~\ref{lem:product-density-bound} therefore gives the bound
\[
 H_0^2:=\exp\left(50\left(\sum_{s\ge1}\kappa_s^2+1\right)\right)<\infty.
\]
Outside the null sets of ambiguous expansions, the mixed-radix map is
a measure-preserving bijection from the reference digit space to $X$.
Transporting the density through this map gives a Haar density with
the same $L^2$ norm. This proves the assertion for both endpoints and
all the deterministic hybrids.
\end{proof}

All the endpoint expansions consequently avoid their ambiguous values
almost surely. Every raw pair satisfies $X_i\le Y_i$, so $A_i\le B_i$.
The nonnegative tail expansion then gives
\begin{equation}\label{eq:coordinate-order}
 u_i(x)\le u_i(y),\qquad
 |u_i(x)-X_i|\le\frac1{r_i^*},\qquad
 |u_i(y)-Y_i|\le\frac1{r_i^*}\quad(i\ge1).
\end{equation}
For $D\ge1$, let $y^{(D)}$ have $y$'s digits at indices $1,\ldots,D$
and $x$'s later digits. Then $u_i(x)\le u_i(y^{(D)})$ for all $i$,
$y^{(D)}\to y$, and $y^{(D)}-x$ is a multiple of $1/q_{D+1}$.

\begin{proposition}[Invariance on the ordered pairs]
\label{prop:ordered-couplings}
Let $p\in M$ be a projection satisfying $(1-p)T_*p=0$. For every $j$,
\begin{equation}\label{eq:coupling-identity}
 (1-p(y))\operatorname{diag}_{k\in\Z}\bigl(R_{x,y}(k)\bigr)p(x)=0
 \quad\text{for }\mathbb P_j\text{-almost every pair}.
\end{equation}
Moreover $0<R_{x,y}(k)\le1$ for all $k\in\Z$ almost surely.
\end{proposition}

\begin{proof}
Fix $D$. The rational shifts $y^{(D)}-x$ have only finitely many possible
values. Lemma~\ref{lem:rational-graphs} holds outside a common Haar-null
set for all these shifts, so the absolute continuity of the $x$ marginal
implies the required identity at $(x,y^{(D)})$. This remains valid even
though the selected shift depends on the pair.

The coordinate order and \eqref{eq:ratio-crossing} bound every ratio by
$1$, both for these pairs and for their limit $(x,y)$. Their measures
converge weakly by the almost-sure convergence $y^{(D)}\to y$.
Lemma~\ref{lem:digit-marginals} supplies uniform $L^2$ bounds for both
marginals. In particular their masses on a set of Haar measure $\delta$
are at most $H_0\sqrt\delta$. Lemma~\ref{lem:compact-transfer} now passes
the identity to the limit.
\end{proof}

We have transferred the invariance identity to ordered pairs. The remaining
reason for choosing these particular laws is that an interior threshold
has a quantitatively positive chance of being crossed, even after one
entire endpoint has been fixed.

\begin{lemma}[Conditional labels and forced crossings]
\label{lem:conditional-labels}
Conditional on $x$, all coordinate labels are independent and each
posterior depends only on the corresponding digit $A_i$. The analogous
statement holds conditional on $y$, using $B_i$.

For $i\in J_s$, $w\in\{a_s,b_s\}$, and $u\in[w,1-w]$, the following
bounds hold almost surely:
\begin{align}
 \mathbb P_j(\text{top label of width }w\text{ at }i\mid x)
   &\ge\rho_s(w)/10 &&\text{if }u_i(x)<u,
   \label{eq:conditional-crossing}\\
 \mathbb P_j(\text{bottom label of width }w\text{ at }i\mid y)
   &\ge\rho_s(w)/10 &&\text{if }u_i(y)\ge u.
   \label{eq:conditional-crossing-bottom}
\end{align}
In the first case that label forces $u_i(y)\ge u$; in the second it
forces $u_i(x)<u$.
\end{lemma}

\begin{proof}
The entire endpoint determines all its digits almost surely. For any
finite collection of labels, conditioning on all $x$-digits gives the
product of the individual posteriors
\[
 \frac{\mathbb P_j(\text{label }\ell\text{ at }i)\,
       \mathbb P_j(A_i=a\mid\text{label }\ell\text{ at }i)}
      {\mathbb P_j(A_i=a)}.
\]
One verifies this formula first against cylinder tests in the digit
sequence, using independence of coordinate pairs, and then against all
bounded digit-measurable tests by the monotone class theorem. The common
initial digit is independent of those pairs. This proves conditional
independence given $x$; the proof for $y$ is the same. In particular,
the dependence of $u_i(x)$ on later digits does not change the label
posteriors after the entire endpoint has been conditioned on.

The parameters satisfy
\[
 \frac1{r_i^*}\le\frac1{N_s}<\frac w2.
\]
If $u_i(x)<u\le1-w$, the whole radix bin containing $X_i$ lies in
$[0,1-w/2)$. On that bin the top component's $X_i$-density is at least
$1$, while the total raw marginal density is at most $10$. Its
posterior probability is therefore at least $\rho_s(w)/10$. Under this
label, \eqref{eq:coordinate-order} gives
\[
 u_i(y)\ge1-w/2-1/r_i^*>1-w\ge u.
\]
For the second assertion, $u_i(y)\ge u\ge w$ places the entire $Y_i$
bin in $[w/2,1)$. The bottom component again has density at least $1$
there, so its posterior is at least $\rho_s(w)/10$. Its $X_i$ value
satisfies
\[
 u_i(x)\le w/2+1/r_i^*<w\le u.
\]
These estimates include thresholds at radix-bin endpoints.
\end{proof}
